\subsection{Time Lax matrix from the finite-lattice construction}
\label{sec:class5-time-lax}

We now take the long-wavelength limit of the finite-lattice time Lax
operator.  The starting point is the operator $A_{a,n}(u)$ defined in
Eq.~\eqref{eq:finite-lattice-time-Lax}.  Its addition by a scalar function of
the quantum spectral parameter leaves the lattice zero-curvature equation
unchanged.  We use the representative
\begin{equation}
\widehat A_{5;a,n}
:=
A_{5;a,n}+\frac{\ii}{u}\Id,
\label{eq:class5-normalized-time-lax}
\end{equation}
where the identity acts on the auxiliary space and on the two physical sites
$n-1,n$ supporting $A_{5;a,n}$.  In terms of the normalized spatial operator
$\widehat L_{5;a,n}$ in Eq.~\eqref{eq:class5-normalized-L},
\begin{equation}
\widehat A_{5;a,n}
=
-\ii\widehat L_{5;a,n}^{-1}
\left(
[h_{5;n-1,n},\widehat L_{5;a,n}]
+\partial_u\widehat L_{5;a,n}
\right).
\label{eq:class5-time-lax-from-Lhat}
\end{equation}
The derivative with respect to $u$ is taken at fixed lattice deformation
parameter $\kappa$.  After imposing the continuum scalings
\eqref{eq:class5-deformation-scaling} and
\eqref{eq:class5-spectral-scaling},
\begin{equation}
\partial_u\widehat L_{5;a,n}
=-\frac{P_{a,n}}{2u^2}
=-\epsilon^2\frac{P_{a,n}}{2\lambda^2}.
\label{eq:class5-u-derivative-Lhat}
\end{equation}

The Hamiltonian on the neighbouring physical sites has the expansion
\begin{equation}
h_{5;n-1,n}
=h^{(0)}_{n-1,n}+\epsilon h^{(1)}_{n-1,n},
\qquad
h^{(0)}:=2P,
\qquad
h^{(1)}:=2\alpha K.
\label{eq:class5-h-expansion}
\end{equation}
Together with the exact relation
$\widehat L_{5;a,n}=\Id+\epsilon L^{(1)}_{5;a,n}$ in
Eq.~\eqref{eq:class5-L-expansion}, Eq.~\eqref{eq:class5-time-lax-from-Lhat}
gives
\begin{equation}
\widehat A_{5;a,n}
=
\epsilon A^{(1)}_{a,n}
+\epsilon^2 A^{(2)}_{a,n}
+\mathcal O(\epsilon^3),
\label{eq:class5-A-expansion}
\end{equation}
with
\begin{align}
A^{(1)}_{a,n}
&=-\ii[h^{(0)}_{n-1,n},L^{(1)}_{5;a,n}],
\label{eq:class5-A1}\\
A^{(2)}_{a,n}
&=
\ii L^{(1)}_{5;a,n}[h^{(0)}_{n-1,n},L^{(1)}_{5;a,n}]
-\ii[h^{(1)}_{n-1,n},L^{(1)}_{5;a,n}]
+\frac{\ii}{2\lambda^2}P_{a,n}.
\label{eq:class5-A2}
\end{align}
Both coefficients act on the tensor product of the auxiliary space and the
two neighbouring physical sites.

We first evaluate the lower symbol of Eq.~\eqref{eq:class5-A-expansion}.
Let $\boldsymbol S_{n-1}$ and $\boldsymbol S_n$ be the coherent-state spin
vectors on the sites $n-1$ and $n$.  Direct contraction of
Eq.~\eqref{eq:class5-A1} gives
\begin{equation}
\bigl(A^{(1)}_{a,n}\bigr)^\downarrow
=
-2
(\boldsymbol S_{n-1}\times\boldsymbol S_n)^i
\frac{\partial U_5}{\partial S^i}
\bigg|_{\boldsymbol S=\boldsymbol S_n}.
\label{eq:class5-A1-lower-symbol}
\end{equation}
Since
\begin{equation}
\boldsymbol S_{n-1}
=
\boldsymbol S(x_n,t)-\epsilon\,\partial_x\boldsymbol S(x_n,t)
+\mathcal O(\epsilon^2),
\end{equation}
we have
\begin{equation}
\boldsymbol S_{n-1}\times\boldsymbol S_n
=
\epsilon\,\boldsymbol S\times\partial_x\boldsymbol S
+\mathcal O(\epsilon^2).
\label{eq:class5-neighbour-cross-product}
\end{equation}
The term $A^{(2)}_{a,n}$ already carries the prefactor $\epsilon^2$ in
Eq.~\eqref{eq:class5-A-expansion}; its lower symbol is therefore needed only
at coincident spin arguments at this order.  Using the lower symbols in
Eqs.~\eqref{eq:class5-P-lower-symbol} and
\eqref{eq:class5-K-lower-symbol}, together with
$\boldsymbol S^{\mathsf T}\boldsymbol S=1$, gives
\begin{equation}
\left.\bigl(A^{(2)}_{a,n}\bigr)^\downarrow\right|_{\boldsymbol S_{n-1}=\boldsymbol S_n=\boldsymbol S}
=
4\ii U_5^2-\frac{\ii}{2\lambda^2}\Id_{\mathcal V_a}.
\label{eq:class5-A2-lower-symbol}
\end{equation}
Equations~\eqref{eq:class5-A1-lower-symbol}--\eqref{eq:class5-A2-lower-symbol}
then yield
\begin{equation}
\widehat A_{5;a,n}^{\downarrow}
=
\epsilon^2 V_5^{\downarrow}(x_n,t;\lambda)
+\mathcal O(\epsilon^3),
\label{eq:class5-A-lower-expansion}
\end{equation}
where
\begin{equation}
V_5^{\downarrow}
=
-2
(\boldsymbol S\times\partial_x\boldsymbol S)^i
\frac{\partial U_5}{\partial S^i}
+4\ii U_5^2
-\frac{\ii}{2\lambda^2}\Id_{\mathcal V_a}.
\label{eq:class5-Vdown-direct}
\end{equation}
This is the continuum matrix obtained from the lower symbol of the finite-lattice
time Lax operator itself.

The remaining contribution comes from the overlapping operator products in the
discrete zero-curvature equation.  Substituting
Eqs.~\eqref{eq:class5-A-expansion} and \eqref{eq:class5-L-expansion} into the
differences defined in Eqs.~\eqref{eq:Rplus-hat-definition} and
\eqref{eq:Rminus-hat-definition}, and keeping the first non-vanishing order,
gives
\begin{align}
&\widehat{\mathcal D}^{(+)}_{a,n}
-\widehat{\mathcal D}^{(-)}_{a,n}
\nonumber\\
&\qquad=
2\ii\epsilon^3\partial_x
\left\{
\left[\bigl(L^{(1)}_5\bigr)^2\right]^\downarrow
-\left[\bigl(L^{(1)}_5\bigr)^\downarrow\right]^2
\right\}
+\mathcal O(\epsilon^4).
\label{eq:class5-product-difference-expansion}
\end{align}
The algebra in Eq.~\eqref{eq:class5-PK-algebra} implies
\begin{equation}
\bigl(L^{(1)}_5\bigr)^2
=
\frac{1}{4\lambda^2}\Id_4,
\label{eq:class5-L1-square}
\end{equation}
and Eq.~\eqref{eq:class5-L-lower-expansion} gives
$\bigl(L^{(1)}_5\bigr)^\downarrow=U_5$.  Hence
\begin{equation}
\left[\bigl(L^{(1)}_5\bigr)^2\right]^\downarrow
-\left[\bigl(L^{(1)}_5\bigr)^\downarrow\right]^2
=
\frac{1}{4\lambda^2}\Id_{\mathcal V_a}-U_5^2,
\label{eq:class5-symbol-square-difference}
\end{equation}
and the continuum contribution defined in
Eq.~\eqref{eq:C-continuum-definition} is
\begin{equation}
\mathcal C_5
=-2\ii\,\partial_x(U_5^2).
\label{eq:class5-shared-site-correction}
\end{equation}

Equation~\eqref{eq:class5-shared-site-correction} is absorbed into the time
component by choosing
\begin{equation}
\Delta V_5
:=
-2\ii U_5^2+\frac{\ii}{4\lambda^2}\Id_{\mathcal V_a}.
\label{eq:class5-DeltaV}
\end{equation}
It obeys
\begin{equation}
[\Delta V_5,U_5]=0,
\qquad
\partial_x\Delta V_5+[\Delta V_5,U_5]=\mathcal C_5.
\label{eq:class5-DeltaV-absorbs-C}
\end{equation}
The scalar term in Eq.~\eqref{eq:class5-DeltaV} fixes the representative; an
$x$-independent scalar function of $\lambda$ leaves the curvature unchanged.
Combining Eqs.~\eqref{eq:class5-Vdown-direct} and
\eqref{eq:class5-DeltaV} gives
\begin{equation}
V_5
=
-2
(\boldsymbol S\times\partial_x\boldsymbol S)^i
\frac{\partial U_5}{\partial S^i}
+2\ii U_5^2
-\frac{\ii}{4\lambda^2}\Id_{\mathcal V_a}.
\label{eq:class5-time-Lax-compact}
\end{equation}
Using the explicit spatial matrix in Eq.~\eqref{eq:class5-spatial-Lax} and the
unit-spin constraint, we write $S_x^a:=\partial_xS^a$ for
$a\in\{+,-,3\}$.  The time Lax matrix is
\begin{equation}
V_5(x,t;\lambda)
=
\begin{pmatrix}
(V_5)_{11}&(V_5)_{12}\\
(V_5)_{21}&(V_5)_{22}
\end{pmatrix},
\label{eq:class5-time-Lax-final}
\end{equation}
with
\begin{align}
(V_5)_{11}
={}&\frac{\ii}{4\lambda^2}
\Bigl[
S^3+\alpha\lambda S^+(1+S^3)+2\alpha^2\lambda^2(S^+)^2
\nonumber\\
&\hspace{3.2em}
+\lambda(S^-S_x^+-S^+S_x^-)
+4\alpha\lambda^2(S^+S_x^3-S^3S_x^+)
\Bigr],
\label{eq:class5-time-Lax-11}\\
(V_5)_{12}
={}&\frac{\ii}{4\lambda^2}
\Bigl[
S^- -\alpha\lambda(1+S^3)^2-2\alpha^2\lambda^2S^+(1+S^3)
\nonumber\\
&\hspace{3.2em}
+2\lambda(S^3S_x^--S^-S_x^3)
+2\alpha\lambda^2(S^+S_x^--S^-S_x^+)
\Bigr],
\label{eq:class5-time-Lax-12}\\
(V_5)_{21}
={}&\frac{\ii}{4\lambda^2}
\Bigl[
S^+ +\alpha\lambda(S^+)^2
+2\lambda(S^+S_x^3-S^3S_x^+)
\Bigr],
\label{eq:class5-time-Lax-21}\\
(V_5)_{22}
={}&\frac{\ii}{4\lambda^2}
\Bigl[
-S^3-\alpha\lambda S^+(1+S^3)
+\lambda(S^+S_x^--S^-S_x^+)
\Bigr].
\label{eq:class5-time-Lax-22}
\end{align}
The difference between Eqs.~\eqref{eq:class5-Vdown-direct} and
\eqref{eq:class5-time-Lax-compact} is the contribution required by the
overlapping operator products in the finite-lattice zero-curvature equation.

